% META: {"id":"1SPE-SUITES-EX-052","chapitre":"1SPE-SUITES","type_objet":"exercice","capacites":["1SPE-SUITES-C3"],"capacites_codes":["C3"],"methodes":[],"parcours":1,"competences":["calculer","raisonner"],"duree_min":6,"mode_creation":"ex_nihilo","sources_inspiration":[],"fichier_tex":"chapitres/1SPE-SUITES/exercices/1SPE-SUITES-EX-052.tex","corrige_tex":"chapitres/1SPE-SUITES/corriges/1SPE-SUITES-CO-052.tex","status":"generated"}
% BEGIN-VERIFY
% # Automatisme : taux d'evolution reciproque.
% from math import isclose
% c = 1.25
% creciproque = 1 / c
% assert isclose(creciproque, 0.8, abs_tol=1e-12)
% assert isclose((creciproque - 1) * 100, -20.0, abs_tol=1e-12)
% # La baisse reciproque n'est pas l'oppose du taux de hausse.
% assert (creciproque - 1) * 100 != -25.0
% # Retour effectif a la valeur initiale.
% assert isclose(200 * c * creciproque, 200, abs_tol=1e-12)
% END-VERIFY
\begin{exercice}{1SPE-SUITES-EX-052}{1}{6}
Le prix d'un article passe de $200$ euros à $250$ euros.
\begin{enumerate}
\item Calculer le taux d'évolution, en pourcentage.
\item Quel taux d'évolution ramènerait le prix de $250$ euros à $200$ euros ?
  On donnera le coefficient multiplicateur, puis le taux en pourcentage.
\item Un client affirme : « Une hausse de $25\,\%$ s'annule par une baisse de
  $25\,\%$. » Justifier que c'est faux, en calculant le prix obtenu.
\item Soit $(u_n)$ la suite géométrique de premier terme $u_0 = 200$ et de
  raison $1{,}25$. Exprimer, à l'aide d'une suite géométrique, le prix obtenu
  après $n$ hausses suivies de $n$ baisses réciproques.
\end{enumerate}
\end{exercice}
